\section{Systematic Uncertainties}%
\label{ch:systematics}

Several sources of systematic uncertainty affect the measured differential
cross-sections. The systematic uncertainties due to detector effects and the
ones related to the modelling  of the signal and background MC components,
which were found to be the most relevant ones, are described in this Section.
Additional MC samples required for some \ttbar and \tW modelling uncertainties
are produced using \textsc{AtlFast-II}. Whenever an \textsc{AtlFast-II} systematic variation sample is used, 
the result is compared to the nominal result derived using the nominal \textsc{AtlFast-II} sample.
Then the systematic uncertainty derived in such comparison is applied to the nominal unfolded result (which uses the full simulation sample).

\subsection{Uncertainty propagation}%
\label{sec:uncertainty_propagation}
Each systematic uncertainty was evaluated at detector-level and then propagated
through the unfolding  procedure (described in \cref{ch:unfolding}).
Deviations from the nominal predictions were evaluated separately for the
upward and downward variations (or in the case of a single variation by
symmetrising the single deviation) for each bin of each observable and 
combining the electron and muon channels.  

In general, the impact on the measurement of each source of uncertainty
is evaluated by unfolding the data with alternative migration matrix, acceptance
and efficiency. The result is then compared with the nominal unfolded distribution.
The same is done for signal modelling systematic uncertainties.

The detector-related uncertainties are briefly described in \cref{sec:syst:detector}
while the uncertainties on the signal and background modelling are discussed in
\cref{sec:syst:modelling,sec:systematics:background}.

\subsection{Detector systematics}%
\label{sec:syst:detector}
The detector systematics can be split into two types: one type where the uncertainty
is assessed by applying event weights to the simulated events (e.g.\ trigger
efficiency uncertainties) and the second type where the uncertainty is assessed
by modifying the four-vector(s) of some of the selected objects  % chktex 36
(e.g.\ jet energy scale uncertainties).

\subsubsection{Leptons}

Electrons and muons both have uncertainties associated with the understanding
of the energy / momentum scale and resolution. These variations modify
the momentum of the leptons. The calibration the lepton momentum reconstruction
and identification comes with associated uncertainties and these are propagated
via changes to the scale factors for the leptons. In a similar way, the scale
factors for the trigger efficiencies are varied to account for uncertainties
in the measurements of the trigger efficiencies of the lepton triggers.

This corresponds to the following nuisance parameters:
\texttt{EG\_RESOLUTION\_ALL},
\texttt{EG\_SCALE\_ALL},\\
\texttt{EL\_EFF\_ID\_TOTAL\_1NPCOR\_PLUS\_UNCOR},
\texttt{EL\_EFF\_Iso\_TOTAL\_1NPCOR\_PLUS\_UNCOR},\\
\texttt{EL\_EFF\_Reco\_TOTAL\_1NPCOR\_PLUS\_UNCOR},
\texttt{EL\_EFF\_TriggerEff\_TOTAL\_1NPCOR\_PLUS\_UNCOR},\\
\texttt{EL\_EFF\_Trigger\_TOTAL\_1NPCOR\_PLUS\_UNCOR},\\
\texttt{MUON\_SAGITTA\_DATASTAT},
\texttt{MUON\_SAGITTA\_RESBIAS},\\
\texttt{MUON\_EFF\_BADMUON\_SYS},
\texttt{MUON\_EFF\_ISO\_STAT},
\texttt{MUON\_EFF\_ISO\_SYS},\\
\texttt{MUON\_EFF\_RECO\_STAT},
\texttt{MUON\_EFF\_RECO\_SYS},
\texttt{MUON\_EFF\_TTVA\_STAT},
\texttt{MUON\_EFF\_TTVA\_SYS},\\
\texttt{MUON\_EFF\_TrigStatUncertainty} and
\texttt{MUON\_EFF\_TrigSystUncertainty}.


\subsubsection{Jets}%
\label{sec:systematics:jets}
The impact of the uncertainty on the jet-energy scale and resolution was
estimated by varying the jet energies according to the uncertainties derived
from simulation and in-situ calibration measurements, following the final
recommendations for the full Run 2 data reprocessed with the ATLAS software
release 21~\cite{JETM-2018-05}.

The uncertainty in the jet energy scale is assessed by using the \texttt{CategoryReduction}
scheme~\cite{web:JetSystematics}, which provides roughly 30 independent variations that
parametrise the uncertainties in the jet energy scale.

This corresponds to the following nuisance parameters:\\
\texttt{JET\_EffectiveNP\_Detector[1-2]},  % chktex 8
\texttt{JET\_EffectiveNP\_Mixed[1-3]},\\  % chktex 8
\texttt{JET\_EffectiveNP\_Modelling[1-4]},  % chktex 8
\texttt{JET\_EffectiveNP\_Statistical[1-6]},\\  % chktex 8
\texttt{JET\_EtaIntercalibration\_Modelling},\\
\texttt{JET\_EtaIntercalibration\_NonClosure\_[2018data, highE, negEta, posEta]},\\
\texttt{JET\_EtaIntercalibration\_TotalStat},\\
\texttt{JET\_Flavor\_Composition\_prop},
\texttt{JET\_Flavor\_Response\_prop},\\
\texttt{JET\_Pileup\_OffsetMu},
\texttt{JET\_Pileup\_OffsetNPV},
\texttt{JET\_Pileup\_PtTerm},\\
\texttt{JET\_Pileup\_RhoTopology} and
\texttt{JET\_PunchThrough\_MC16}.

The uncertainty in the jet energy resolution is assessed using the \texttt{FullJER}
scheme~\cite{web:JetSystematics}. This provides 13 independent variations that parametrise
the uncertainties in the jet energy resolution, smearing both (pseudo-)data  % chktex 36
and MC events.

This corresponds to the following nuisance parameters:\\
\texttt{JET\_JER\_DataVsMC\_MC16} and
\texttt{JET\_JER\_EffectiveNP\_[1-12]}.  % chktex 8

The uncertainty originating from the JVT requirement on jets is assessed by
varying the scale factor associated with the JVT requirement~\cite{ATLAS-CONF-2014-018}.
This corresponds to the following nuisance parameter:
\texttt{JET\_JvtEfficiency}.

\subsubsection{\btag}%
\label{sec:systematics:btag}
The identification of jets which originate from a \bhadron is performed using
a multivariate algorithm, \texttt{DL1r}~\cite{web:FlavourTagging}.
Scale factors are applied to simulation to correct the jet flavour tagging
efficiencies to match those in data~\cite{FTAG-2018-01, ATLAS-CONF-2018-001, ATLAS-CONF-2018-006}.
The corrections are applied for the tagging efficiency of \bjets and
the mis-tag rates of \(c\)- and light-jets, in bins of the pseudo-continuous b-tagging scores,
with corresponding uncertainties derived for each jet flavour.
A principle component analysis is performed to diagonalise the covariance
matrix given the input systematic uncertainties considered in the calibration
of the tagging efficiencies. This results in a number of orthogonal eigen-variations
for each jet flavour, which are then used as nuisance parameters in this analysis.

This corresponds to the following nuisance parameters:\\
\texttt{FT\_EFF\_Eigen\_B\_[0-8]},  % chktex 8
\texttt{FT\_EFF\_Eigen\_C\_[0-3]},  % chktex 8
\texttt{FT\_EFF\_Eigen\_Light\_[0-3]},\\  % chktex 8
\texttt{FT\_EFF\_extrapolation} and
\texttt{FT\_EFF\_extrapolation\_from\_charm}.


\subsubsection{Missing transverse energy}%

The uncertainties on the energy scales and resolutions of the reconstructed
leptons and jets are propagated to the missing transverse momentum
in a correlated way. 

In addition, there are scale and resolution uncertainties associated with the
soft term that are independent uncertainty components~\cite{ATLAS-CONF-2018-023}.
These uncertainties modify the magnitude and direction of the missing transverse
momentum vector.

This corresponds to the following nuisance parameters:\\
\texttt{MET\_SoftTrk\_ResoPara},
\texttt{MET\_SoftTrk\_ResoPerp} and
\texttt{MET\_SoftTrk\_Scale}.
\\The first two are symmetrised.

\subsubsection{Pileup modelling}
The uncertainty on the reweighting procedure used to correct the pile-up profile
in MC to match the data is obtained by varying by \(\pm \qty{1}{\sigma}\)
the scale factors obtained with the \texttt{PileupReweightingTool}
\footnote{\url{https://twiki.cern.ch/twiki/bin/view/AtlasProtected/ExtendedPileupReweighting}}.
This corresponds to the following nuisance parameter: \texttt{PRW\_DATASF}.

\subsection{Signal modelling uncertainties}%
\label{sec:syst:modelling}

Limited theoretical accuracy in the modelling of the signal processes
is a significant source of uncertainty. The impact of this kind of uncertainties
has been evaluated using alternative predictions that are described
in \cref{sec:datamc:mc} following the procedure described in \cref{sec:uncertainty_propagation}. 

\subsubsection{Cross-section}
While in principle the signal is normalised during the unfolding process
two samples are stitched together with their independent normalisations.
To account for possible fraction differences due to cross-section uncertainties
\ttbar and \tW cross-section uncertainties are anti-correlated and this
is applied as an uncertainty.

In practice the following variations are taken and samples are rescaled with:

\begin{itemize}
  \item \ttbar \( +1.03519 \), \tW \( -1.03531 \)
  \item \ttbar \( -1.04391 \), \tW \( +1.03657 \)
\end{itemize}

This corresponds to the following nuisance parameter: \texttt{THEORY\_CROSS\_SECTION\_signal}.

\subsubsection{Hadronisation model}%
\label{sec:systematics:partonshower}
The uncertainty due to the choice of the hadronisation model and other
non-perturbative aspects of the parton shower is evaluated using the
\POWHEG{}+\HERWIG[7] sample.

This corresponds to the following nuisance parameter: \texttt{THEORY\_SHOWERING\_HERWIG7\_signal}.
\\The uncertainty is symmetrised and uses dedicated \textsc{AtlFast-II} samples.

\subsubsection{Scales}%
\label{sec:systematics:scale}
The uncertainty due to the choice of the \muR and \muF scales in
the matrix element calculation is evaluated using an independent variation
of \muR or \muF by a factor 0.5 or 2.0 as up and down variations.

The ISR modelling uses the the Var3c eigentune variation of the A14
tune~\cite{ATL-PHYS-PUB-2017-007}.

The uncertainty due to the FSR simulation is obtained by varying
the renormalization scale used in the final-state shower
by a factor 2.0 and 0.5 with respect to the nominal value.

This corresponds to the following nuisance parameters:
\texttt{THEORY\_SCALE\_FACTORISATION\_signal},
\texttt{THEORY\_SCALE\_RENORMALISATION\_signal},
\texttt{THEORY\_ISR\_signal} and
\texttt{THEORY\_FSR\_signal}.


\subsubsection{\texorpdfstring{\hdamp}{Hdamp}}%
\label{sec:systematics:hdamp}
Another parameter affecting the additional radiation is \hdamp.
The uncertainty due to the choice of this parameter is estimated independently
from the other ISR components using a dedicated \POWHEG{}+\PYTHIA{8} sample
where \hdamp\ is multiplied by a factor 2 with respect to the nominal value.

This corresponds to the following nuisance parameter: \texttt{THEORY\_HDAMP\_signal}.
\\The uncertainty is symmetrised and uses dedicated \textsc{AtlFast-II} samples.

\subsubsection{Matching}%
\label{sec:systematics:matching}
The matching uncertainty is evaluated by comparing the nominal sample with
an alternative sample obtained setting the \pthard parameter in \PYTHIA to 1
and comparing it to the nominal sample, in which the parameter is set to 0.
This parameter regulates the definition of the vetoed region of the showering,
which is important to avoid hole/overlap in the phase space filled by \POWHEG
and \PYTHIA. This prescription is based on the studies described in~\cite{10.21468/SciPostPhys.12.1.010}.

This corresponds to the following nuisance parameter: \texttt{THEORY\_PTHARD\_signal}.
\\The uncertainty is symmetrised and uses dedicated \textsc{AtlFast-II} samples.

% \subsection{Lineshape}%
% \label{sec:systematics:lineshape}
% The lineshape uncertainty of the top pair system is evaluated by comparing
% the nominal sample with an alternative sample in which the top quarks are
% decayed with \madspin.
% This is the modelling uncertainty in which only the top pair component
% of the signal is modified; the nominal \tW sample is used in the alternative
% sample.

% This corresponds to the following nuisance parameter: \texttt{THEORY\_MADSPIN\_signal}.

\subsubsection{Recoil to top}%
\label{sec:systematics:recoil-to-top}
The modelling of the parton shower of the \bquarks in \( t \rightarrow W b \)
is affected by the gluon recoil scheme. The default samples employ the recoiling
against \bquarks. This affects the second and subsequent gluon emission from
quarks produced by coloured resonance decays, such as the \bquark in
a \( t \rightarrow W b \) process, deciding how the momentum is re-arranged
between the \Wboson and the \(b\). This uncertainty compares the nominal sample
with the sample where the top-quark itself works as recoiling object for the
second and subsequent emissions. This option give wider-angle gluon radiation,
resulting in energy deposits that do not get clustered into the \bjet,
and lower gluon-energy emission, altering the modelling of the \bquark
fragmentation and hardening the \bhadron momenta.
This is the modelling uncertainty in which only the top pair component
of the signal is modified; the nominal \tW sample is used in the alternative
sample.

This corresponds to the following nuisance parameter: \texttt{THEORY\_TOPRECOIL\_signal}.
\\The uncertainty is symmetrised and uses dedicated \textsc{AtlFast-II} samples.

\subsubsection{Top-quark mass}%
\label{sec:systematics:topmass}
The top-quark mass is known with a precision of around \qty{0.5}{\GeV}~\cite{TOPQ-2017-03}.
This can impact the analysis because the kinematics of the decay products
(and hence the probability to pass the selection requirements) depends
on the top-quark mass. The uncertainty is estimated by using two additional
simulated samples with varied top-quark mass values of \qty{169}{\GeV} and
\qty{176}{\GeV}. These samples otherwise have the same setup as the nominal
\ttbar and \tW samples and are used to evaluate an uncertainty via
the unfolding result, identically to other modelling effects.
The final uncertainty is then scaled by 1/7 to match the effect of
a \( \pm \qty{0.5}{\GeV} \) shift in the top-quark mass.
The uncertainty of \qty{0.5}{\GeV} is obtained from events in the full phase-space.
As the top-quark mass is Lorentz invariant, it is appropriate to apply this
top-quark mass value and uncertainty to the limited phase-space used in the analysis. 

This corresponds to the following nuisance parameter: \texttt{THEORY\_MTOP\_signal}.
\\The uncertainty uses dedicated \textsc{AtlFast-II} samples.

\subsubsection{PDF}%
\label{sec:systematics:pdf}
The effect of a different PDF choice modifies the efficiency, acceptance and
potentially also the response matrix, i.e.\ the corrections used to correct
the spectrum at the detector level to the particle level. 
The impact of the choice of different PDF sets has been assessed by applying
an event-by-event reweighting procedure to the nominal \POWHEG{}+\PYTHIA{8}
\ttbar and \tW samples following the \PDFforLHC[15] prescription~\cite{Butterworth:2015oua}
using the \texttt{LHAPDF} tool~\cite{Buckley:2014ana}.

The variations are applied as weight variations separately for PDF and separately for \alphas.
The ensemble of histograms is fed to the \texttt{LHAPDF} tool together with the relevant
PDF set used at the generation. The tool then yields combined PDF+\alphas uncertainty.

This corresponds to the following nuisance parameter: \texttt{THEORY\_PDF4LHC\_VARIATION\_signal}.

\subsubsection{\tW DR vs DS}%
\label{sec:systematics:DRvsDS}
Another component of the uncertainty is represented by the difference between
the two approaches diagram removal and diagram subtraction, used to remove
the overlap between the \ttbar and \tW productions. This uncertainty only
affects the \tW component of the signal; the difference between the two
\tW DR/DS samples is used to form the uncertainty estimate, with \ttbar
remaining fixed to the nominal value.

This corresponds to the following nuisance parameter: \texttt{THEORY\_DR\_DS\_signal}.
\\The uncertainty is symmetrised.

\subsubsection{Heavy flavour uncertainty}
The uncertainty estimates the effect of mismodelling of \ttbar \(+ b\) and
\ttbar \(+ c\) events which may be underestimated in MC by scaling up
such events by \qty{30}{\%} and \qty{60}{\%}, respectively.
As the uncertainty is defined a bit ad-hoc and the overall effect is
very small, it is not applied by default.
See more details in \cref{app:heavyFlavour}.

\subsection{Uncertainties on the background processes}%
\label{sec:systematics:background}

Several backgrounds enter the measurement region estimation, therefore several
uncertainties are associated to them.

\subsubsection{\texorpdfstring{\Wjets}{W+jets}}

Full detector uncertainties and the modelling uncertainties on PDF, scales
and electroweak correction are applied similar to the signal sample.

A likelihood fit in a dedicated control region is then performed to
constrain the normalisation effects of the individual uncertainties
and splitting the shape and normalisation effects.
Only shape components of all uncertainties on \Wjets are taken into account.
Additionally a nuisance parameter is added to account for post-fit uncertainty
on the \Wjets normalisation.
More details are described in \cref{sec:LikelihoodFit}.

Electroweak correction uncertainty is estimated by taking the envelope
of the additive, multiplicative and exponentiated approach of calculating
the uncertainty, where the exponentiated is the nominal way to apply
the correction.

This corresponds to the following additional nuisance parameters:

\paragraph{THEORY\_SCALE\_COMBINED\_Wjets} 7-point QCD scale variations stored as internal weights in the baseline samples, from which the envelope around the nominal can be worked out. The scale uncertainty is then given by the maximum shift of the envelope with respect to the nominal. 

\paragraph{THEORY\_PDF4LHC\_VARIATION\_Wjets} Parton Density Function The Hessian NNPDF3.0 NNLO uncertainty scheme is to be adopted, which can be accessed via the LHAPDF tool as it has built-in methods to work out the combined PDF uncertainty~\cite{Buckley:2014ana}. The central PDF values for the NNPDF3.0nnlo (PDF261000), MMHT2014nnlo68cl (PDF25300) and CT14nnlo (PDF13000) PDF sets are included as variations as well as two NNPDF3.0nnlo alpha\_s variations (PDF269000 and PDF270000).

\paragraph{THEORY\_EWK\_Wjets} Electroweak Virtual Corrections: Approximate NLO EW corrections are included as on-the-fly weights using the electroweak virtual approximation.


\subsubsection{Other background processes}

Other background processes overall contribute less than \qty{5}{\%} of the total
yield in the regions and do not have a big effect on the total uncertainty.
The ``\OtherV'' samples have scales, PDF and electroweak
uncertainties applied, while the ``\OtherT'' samples just have
a conservative \qty{25}{\%} normalisation uncertainty applied.

This corresponds to the following additional nuisance parameters:\\
\texttt{THEORY\_SCALE\_COMBINED\_multiboson\_noW},
\texttt{THEORY\_PDF4LHC\_VARIATION\_multiboson\_noW},\\
\texttt{THEORY\_EWK\_multiboson\_noW} and 
\texttt{THEORY\_CROSS\_SECTION\_other\_top\_noWt}.

\subsubsection{Background from fake and non-prompt leptons}%
\label{sec:Systematics_fakes}

Two uncertainties are assigned to the fake and non-prompt lepton background,
separately for electrons and muons. They are derived based on variations
in fake estimation selection and described more in detail in
\cref{sec:bkg-fakes-electron-sys,sec:bkg-fakes-muon-sys}.

This corresponds to the following nuisance parameters:
\texttt{FAKES\_Electron} and 
\texttt{FAKES\_Muon}.

\subsection{Other uncertainties}

\subsubsection{Finite MC sample statistics}%
\label{sec:systematics:finiteMC}
To account for the limited statistics of the Monte Carlo samples, pseudo-experiments are used to evaluate the impact of finite statistics. This uncertainty is evaluated on both signal and background samples. 
The number of events in each bin is smeared by a Gaussian shift with mean equal to the yield of the bin, and  standard deviation equal to the uncertainty of the bin. Then, the smeared spectrum is unfolded.
The procedure is replicated \(10\)k times, then the final statistical uncertainty is evaluated from the quantiles over all the toy distributions. 
The resulting statistical uncertainty was found to be typically below \qty{0.5}{\%}, increasing to \(1\)-\qty{2}{\%} in the tails of some distributions.\@ \cref{fig:StatUncMC:ptl1_postUnfold} shows the statistical uncertainty as a function of lepton \pt.

\begin{figure}[htbp]
  \centering
  \includegraphics[page=13, width=0.45\textwidth]{BulkSelection/TUnfoldStandalone_OptionA_data_nonClosureAlternative_SR_Bulk_Final_l_4j_incl_STAT_MC}
  \caption{MC statistical uncertainty for the lepton \pt variable in the bulk region.}%
  \label{fig:StatUncMC:ptl1_postUnfold}
\end{figure}

% chktex 8
\subsubsection{Integrated luminosity}%
The uncertainty in the combined 2015--2018 integrated luminosity is \qty{0.82}{\%}~\cite{DAPR-2021-01},
obtained using the LUCID-2 detector~\cite{LUCID2} for the primary luminosity measurements,
complemented by measurements using the inner detector and calorimeters.
This uncertainty is propagated to the unfolded results by varying the luminosity
by \(\pm \qty{0.82}{\%}\) in the reconstructed simulated spectra and unfolding
them using the nominal corrections.

\subsubsection{Data statistics}%
\label{sec:data_statistics}

The uncertainty due to the data statistics is propagated through the unfolding by creating Poisson smeared pseudo-experiments \(N_\mathrm{pseudo-exp} = 10000\) that are passed through the  % chktex 35
unfolding procedure.

The shifting of data is defined as:

\begin{equation}
    X_\mathrm{shift} = X - \mathrm{Pois}(X)
\end{equation}

\(X\) is the number of events in the bin. To evaluate the data statistical uncertainty
it is shifted by \( \mathrm{Pois}(X) \), a random number from the Poisson
distribution with the mean \(X\).

\cref{fig:StatUncData:ptl1_postUnfold} shows the data statistical uncertainty as a function of lepton \pt. 
The resulting statistical uncertainty is found, for this variable, to be below \qty{4}{\%}.

\begin{figure}[htbp]
  \centering
  \includegraphics[page=13, width=0.45\textwidth]{BulkSelection/TUnfoldStandalone_OptionA_data_nonClosureAlternative_SR_Bulk_Final_l_4j_incl_STAT_DATA}
  \caption{Data statistical uncertainty for the lepton \pt variable in the bulk region.}%
  \label{fig:StatUncData:ptl1_postUnfold}
\end{figure}


\subsubsection{Mass bias in the hadronic \texorpdfstring{\Wboson}{W} reconstruction}%
\label{sec:WMassVar}
The uncertainty discussed in this section is related to the reconstruction algorithm introduced in \cref{sec:WhadReco}.
Therefore, it is only considered in the \Wboson-hadronic measurement region defined in \cref{sec:regions-wbwb}.
An explicit assumption on the mass of the \Wboson boson is used in the reconstruction algorithm.
However, the \Wboson mass cannot be identified with a fixed value due to its natural width and resolution effects.
This study investigates the effect of a shifted mass assumption in the algorithm.
The value of the mass is shifted by \(\qty{4}{\GeV} \approx 2\Gamma_{\Wboson}\) up or down.
The change of the detector level distributions is propagated to the unfolded cross section values.
The resulting change in the cross section is displayed in \cref{fig:WMassVar} as a function of \Wboson boson \pt and \(\pt(b\Wboson)\).
The uncertainty has the largest impact on these two observables. It is still well below \qty{1}{\%} in all bins.
It is negligible for all other measured observables.
\begin{figure}[h]
  \centering
  \subfloat[]{
    \includegraphics[page=70, width=0.4\textwidth]{figures/Systematics/TUnfoldStandalone_OptionA_data_nonClosureAlternative_SR_Whad_Final_l_Whad_particle_WMASS_VAR_signal}%
    \label{fig:WMassVar:Wpt_postunfold}
  }
  \subfloat[]{
    \includegraphics[page=124, width=0.4\textwidth]{figures/Systematics/TUnfoldStandalone_OptionA_data_nonClosureAlternative_SR_Whad_Final_l_Whad_particle_WMASS_VAR_signal}%
    \label{fig:WMassVar:bWpt_postunfold}
  }
  \caption{Uncertainty resulting from varying the mass in the hadronic \Wboson reconstruction as a function of 
  \protect\subref{fig:WMassVar:Wpt_postunfold} \(\pt(\Wboson)\) and \protect\subref{fig:WMassVar:bWpt_postunfold} 
  \(\pt(b\Wboson)\).}%
  \label{fig:WMassVar}
\end{figure}
